% ===========================================================================
\chapter{The Learning Rule Decides How Many Units Carry the Solution}
\label{chap:grokking}
\margintoc
% ===========================================================================

%% TODO(story): written by the framing pass. Placeholder until then.
\begin{publicationbasis}{Preprint}
Full citation of the paper this chapter is built on, ending in its
\texttt{\textbackslash cite} key. See \texttt{imports/grokking.toml}.
\end{publicationbasis}

%% TODO(story): written by the framing pass. Placeholder until then.
\contributionstatement

Placeholder. Who did what, factually, naming co-authors.

%% TODO(story): written by the framing pass. Placeholder until then.
\section{Chapter Introduction}
%% TODO(story): this section has to define the twelve rules the chapter tests,
%% and say which eight vary the optimizer. The paper defines them in the last
%% subsection of its background section, which the chapter drops as covered by
%% the thesis background chapter, and the paper's body refers back to that
%% definition as Section~\ref{sec:ruleset}. The label is placed here because
%% this is where the definition now lives. If the definition ends up in the
%% background chapter instead, move the label with it.
\importedlabel{grokking}{sec:ruleset}

Placeholder. Connects the previous chapter's result to this one and
states what the chapter contributes to the thesis argument.

% ---------------------------------------------------------------------------
% The published paper, included verbatim. Do not edit imported_papers/.
% ---------------------------------------------------------------------------

\begin{importedpaper}{grokking}
\importpart{03_setup}
\importpart{04_whether-a-rule-groks-is-decided-at-the-minima-of-the-loss}
\importpart{05_the-rules-that-grok-share-the-algorithm-but-not-the-circuit}
\importpart{06_the-two-ends-trade-robustness-against-adaptation-speed}
\end{importedpaper}

%% TODO(story): written by the framing pass. Placeholder until then.
\section{Chapter Summary}

Placeholder. States the chapter's findings once in full and hands off
to the next chapter.

